\documentclass[11pt]{article}
\usepackage[margin=1in]{geometry}
\usepackage{amsmath,amssymb,amsthm,mathtools,booktabs,hyperref}
\newtheorem{definition}{Definition}
\newtheorem{theorem}{Theorem}
\newtheorem{lemma}{Lemma}
\newtheorem{proposition}{Proposition}
\newtheorem{remark}{Remark}
\title{Robust and Epistemic Viability for Hybrid Material--Institution Systems}
\author{Author(s) redacted for review}
\date{}
\begin{document}\maketitle

\begin{abstract}
This paper develops conditional mathematical kernels for hybrid sustainability systems with material histories, incomplete observations, institutional authority, and uncertain implementation. The central state for decision is an information state consisting of compatible physical histories and parameters together with an institutional mode. We establish boundary-relative hybrid moiety balance and conditional hybrid positivity, define sampled epistemic-institutional predecessor operators, and prove a fixed-point characterization of finite- and infinite-horizon viability under explicitly stated transition and strategy assumptions. We prove information-refinement monotonicity and a common-action obstruction that distinguishes physical, epistemic, authority, and implementation infeasibility. We also state conditional observer-to-safety and compositional-safety results. Theorems are conditional: construction of compatible-state updates, verification of nonempty kernels, and calibration of a domain module remain empirical and model-specific obligations.
\end{abstract}

\section{Status and setting}
This paper is the mathematical companion to Paper I. It does not claim that every sustainability system has a nonempty epistemic-institutional viability kernel. It supplies a conditional formal object in which material feasibility, information, institutional authority, implementation, and safety are jointly represented.

Let the physical history be $\phi_t=X_t\in\mathscr H=C([-\tau_{\max},0],\mathcal X)$, where $X=(r,f,h)$ follows the hybrid architecture of Paper I. An information state is
\[
Z=(B,h)\in\mathfrak S,
\qquad B\subseteq\mathscr H\times\Theta,
\]
where $B$ contains all compatible physical histories and parameters. Decision times $t_k$ are locally finite.

\section{Hybrid material feasibility}
Between events, material dynamics satisfy
\[
\dot r=\mathsf S\nu(\phi_t,u,\omega,\vartheta)+b(\phi_t,u,\omega,t),
\]
with $\nu\ge0$, separate reverse columns, and donor-limited negative boundary flows. At events $r^+=\Delta_r(\cdot)$.

\begin{theorem}[Conditional hybrid moiety balance]
Suppose $r$ is absolutely continuous between locally finite event times and has left and right limits at events. If $\mathsf L^\top\mathsf S=0$, then
\[
\mathsf L^\top r(t)-\mathsf L^\top r(0)
=
\int_0^t\mathsf L^\top b(\phi_s,u(s),\omega(s),s)\,ds
+\sum_{t_k\le t}\mathsf L^\top[r(t_k^+)-r(t_k^-)].
\]
\end{theorem}
\begin{proof}
Integrate the continuous balance between consecutive events and telescope the left/right state differences.
\end{proof}

\begin{theorem}[Conditional hybrid history-cone invariance]
Assume: (i) the RFDE flow is well posed and unique between locally finite event times; (ii) every nonnegative history with $\phi_j(0)=0$ satisfies the material quasipositivity condition $[\mathsf S\nu(\phi,\ldots)+b(\phi,\ldots)]_j\ge0$; and (iii) every admissible reset map sends nonnegative material pre-event states to nonnegative post-event states. Then nonnegative initial material history generates a nonnegative hybrid material trajectory on its interval of existence.
\end{theorem}
\begin{proof}
The history-cone condition prevents continuous first exit between events; reset positivity prevents exit at events. Induction over the locally finite event sequence gives the result.
\end{proof}

\begin{remark}
Positivity, boundedness, persistence above a threshold, and viability are distinct. The theorem establishes only positivity.
\end{remark}

\paragraph{Yield-routing obligation.} If a transformation is represented with a yield below one for a declared moiety, the omitted fraction must be routed to another represented compartment or a declared boundary flow. Otherwise the claimed moiety balance applies only after silently dropping that moiety from $\mathsf L$.

\section{Ideal full-information benchmark}
For a finite-dimensional fully observed ODE module $\dot X=f(X,u,w)$ with state-safe set $\mathcal C_X$, let $T_{\mathcal C_X}(X)$ denote its contingent cone.

\begin{proposition}[Conditional full-information viability benchmark]
Suppose $\mathcal C_X$ is closed and sufficiently regular for the stated tangent representation, or suppose the direct condition $f(X,u,w)\in T_{\mathcal C_X}(X)$ is verified. If at every boundary state there exists a common admissible control for all active constraints and all $w\in\mathcal W$, the safe-control correspondence admits an appropriate regular selection or viable solution concept, and solutions exist on $[0,T]$, then standard controlled-invariance results establish full-information viability on that horizon.
\end{proposition}
\begin{remark}
This is an ideal ODE benchmark. It is not a theorem for delayed, sampled, partially observed institutions.
\end{remark}

\begin{lemma}[Stability and safety are independent]
Local or asymptotic stability of an equilibrium does not imply that it belongs to a declared safe set, and membership of a nominal equilibrium in a safe set does not imply robust safety under uncertainty.
\end{lemma}
\begin{proof}
Take $\dot x=-(x+1)$ with safe set $[0,\infty)$ for the first statement. For the second, take a nominally safe stable system and permit a disturbance class that drives the state outside the set.
\end{proof}

\section{Institutional game and quantifier convention}
A prescription and an implemented control are distinct:
\[
a\in\Gamma(B,h),\qquad u\in\mathcal E(B,h,a).
\]
Here $\Gamma$ is institutional authority and $\mathcal E$ is the implementation/compliance correspondence. Robust safety always uses the lower-game order
\[
\exists\text{ non-anticipative prescription strategy}
\quad\forall\text{ admissible implementations, disturbances, parameters, and observation branches}.
\]
The reverse order $\forall w\exists u_w$ is not an implementable claim unless $w$ is observed before action.

Let $\operatorname{TubeSafe}(B,h,a)$ mean that every compatible trajectory under every $u\in\mathcal E(B,h,a)$ remains in the physical safe set and satisfies state-control admissibility until the next event. Let $\Psi(B,h,a,Y^+)$ be the compatible-state update after the next observation and $\Delta_h$ the institutional update.

\section{Sampled epistemic-institutional predecessor}
For $\mathfrak Q\subseteq\mathfrak S$, define
\[
\operatorname{Pre}_{\mathfrak I}(\mathfrak Q)=
\left\{(B,h):\begin{array}{l}
\exists a\in\Gamma(B,h)\text{ such that }\operatorname{TubeSafe}(B,h,a),\\
(\Psi(B,h,a,Y^+),\Delta_h(h,a,Y^+,\rho))\in\mathfrak Q\\
\text{for every compatible observation, covariate, and implementation outcome}
\end{array}\right\}.
\]

\begin{definition}[Finite-horizon epistemic-institutional kernel]
Let $\mathfrak K_0=\mathfrak S$ and use the safe-base recursion
\[
\mathfrak K_{n+1}=\mathcal T(\mathfrak K_n):=\mathfrak S\cap\operatorname{Pre}_{\mathfrak I}(\mathfrak K_n).
\]
The deflationary implementation $\mathfrak K_n\cap\operatorname{Pre}_{\mathfrak I}(\mathfrak K_n)$ generates the same descending orbit from $\mathfrak S$ when $\operatorname{Pre}_{\mathfrak I}$ is monotone; this follows inductively because each later predecessor is contained in the predecessor of the preceding larger iterate. The safe-base form is official because it remains safe under arbitrary initialization.
\end{definition}

\begin{theorem}[Conditional sampled epistemic-institutional viability]
Assume: (i) the compatible-state update is nonempty for every compatible branch; (ii) all transition and authority/implementation correspondences are defined on $\mathfrak S$; (iii) the required non-anticipative prescriptions exist whenever the predecessor condition holds; and (iv) $\operatorname{Pre}_{\mathfrak I}$ is monotone. Then $\mathfrak K_n$ is the set of information states from which safety can be guaranteed for $n$ decision intervals. Moreover the displayed monotone operator $\mathcal T$ has a greatest fixed point on the powerset lattice of $\mathfrak S$.

If, in addition, $\mathcal T$ preserves decreasing countable intersections (or an equivalent $\omega$-continuity-from-above condition) and the strategy-selection assumptions are closed under the infinite-horizon limit, then
\[
\mathfrak K_\infty:=\bigcap_{n\ge0}\mathfrak K_n
\]
is that greatest fixed point. Without this additional condition, the greatest fixed point is obtained by the standard transfinite descending iteration rather than necessarily by the countable intersection displayed above.
\end{theorem}
\begin{proof}
Induct on $n$: membership in $\mathfrak K_{n+1}$ supplies a current safe prescription and a successor in $\mathfrak K_n$ for every branch. Monotonicity makes $\mathcal T$ monotone, so Tarski's theorem gives a greatest fixed point. The additional continuity and closure hypotheses identify the countable descending limit with that fixed point; otherwise transfinite iteration is required.
\end{proof}

\begin{remark}
The theorem is an abstract characterization, not a claim that $B_t$ is tractable, that $\mathfrak K_\infty$ is nonempty, or that a particular ecological system meets the assumptions.
\end{remark}

\section{Information, common action, and observer margins}
\begin{theorem}[Information-refinement monotonicity]
Suppose information structure $\mathcal I_1$ causally refines $\mathcal I_2$, while physical dynamics, authority, implementation correspondence, uncertainty class, and safety constraints are unchanged. If a controller may ignore additional information, then
\[
\operatorname{IViab}^{\mathcal I_2}_T\subseteq\operatorname{IViab}^{\mathcal I_1}_T.
\]
\end{theorem}
\begin{proof}
Every $\mathcal I_2$-measurable strategy is also an $\mathcal I_1$-measurable strategy by ignoring the additional information.
\end{proof}

For an information set $B$, let $\mathcal U_{\rm com}(B)$ be the intersection of all robustly safe implemented actions across compatible states, subject to the next-review tube condition.

\begin{proposition}[Common-action obstruction]
Suppose no informative observation arrives before the next action must be selected. If $\mathcal U_{\rm com}(B)=\varnothing$, then $(B,h)$ is not robustly viable under output feedback, even if every compatible physical state is individually viable under full information.
\end{proposition}
\begin{proof}
Output feedback must choose one action/prescription before the uncertainty within $B$ is resolved. No such action is robustly safe for all compatible states.
\end{proof}

\begin{proposition}[Conditional observer-to-safety transfer]
Let a full-state feedback law have a uniform inward margin $\eta_i>0$ for active safety constraint $b_i$. Suppose the implementation error due to an estimator obeys
\[
|\nabla b_i(X)\,[f(X,k(\widehat X),w)-f(X,k(X),w)]|\le L_i\|\widehat X-X\|.
\]
Then the inward inequality is preserved whenever $L_i\|\widehat X-X\|\le\eta_i$.
\end{proposition}
\begin{proof}
Subtract the estimation-induced error bound from the full-state inward margin.
\end{proof}
\begin{remark}
This is a local sufficient condition. It motivates eroded safe sets but does not supply an observer or establish estimator bounds.
\end{remark}

\section{Information value, recovery, compositionality, and discrepancy}
For a noncompensatory safety margin $q(X,u)=\min_i m_i(X,u)$, define the robust information value, when the optimization is well posed, by
\[
V_T^{\mathcal I}(B,h)=\sup_{\Pi\in\Pi_{\mathcal I}}\inf_{\rm compatible\ branches}\min_{0\le t\le T}q(X(t),u(t)).
\]
A nonnegative value is equivalent to finite-horizon robust safety only when the supremum is attained or the relevant viability closure is used. For an information refinement, the difference $V_T^{\mathcal I_1}-V_T^{\mathcal I_2}$ is a safety-margin value of information, not an entropy metric.

\paragraph{Recovery template.} Let $\mathcal E_X\supseteq\mathcal C_X$ be an emergency envelope. Informational recovery asks whether a strategy keeps the physical trajectory in $\mathcal E_X$ until some time $T$ and reaches a viable information state $(B_T,h_T)\in\mathfrak K_\infty$. Physical recoverability therefore need not imply institutional recoverability.

\paragraph{Safe-learning template.} Compatible-state updates depend on action. An action is safely informative only if it is tube-safe and contracts a declared belief-size functional for every compatible observation branch. This is a domain-specific dual-control problem; learning is not presumed harmless.

\begin{proposition}[Conditional compositional safety]
Suppose each subsystem has a robust invariant certificate under specified interface assumptions. If all material interfaces are consistently signed and routed, every subsystem guarantee satisfies the assumptions of connected subsystems, shared controls are jointly feasible, and hybrid event composition is nonblocking, then the product certificate is invariant for the interconnected system.
\end{proposition}
\begin{remark}
The proposition is conditional on the interface contracts; separate subsystem certificates alone do not imply network safety.
\end{remark}

\paragraph{Structural discrepancy template.} Domain modules may use
\[
\dot X=F(X_t,u,w,\vartheta)+\delta_t,
\qquad \delta_t\in\mathcal D_t,
\]
but must separately report process uncertainty, parameter uncertainty, observation error, structural discrepancy, and implementation uncertainty. Shrinking $\mathcal D_t$ requires a stated coverage argument, not fitted-model confidence alone.

\paragraph{Normative monotonicity.} If two declared state/control constraint systems satisfy $\mathcal C_X(\lambda_1)\subseteq\mathcal C_X(\lambda_2)$ and $\mathcal A(\cdot;\lambda_1)\subseteq\mathcal A(\cdot;\lambda_2)$, then the corresponding viability kernels are nested in the same direction under aligned information, authority, and uncertainty classes. This exposes feasibility consequences of normative choices without purporting to derive them.

\section{Delay, reduction, and nonlinear transitions}
\begin{proposition}[No sign-free delay conclusion]
Delay magnitude alone has no universal stabilizing or destabilizing implication. The scalar systems $\dot x=-ax-bx(t-\tau)$ and $\dot x=-ax+bx(t-\tau)$, with $a,b>0$, have different feedback signs and different stability properties. Any policy claim additionally depends on plant dynamics, gain, sampling, implementation, and constraints.
\end{proposition}
\begin{remark}
This is a non-universality statement, not an analysis of a particular sampled institution, distributed delay, non-minimum-phase plant, or ecological system.
\end{remark}

\paragraph{Conditional ODE reduction result.} Classical finite-horizon Tikhonov reduction may be invoked for a specified ODE slow-fast subsystem only after its isolated fast equilibrium, uniform Hurwitz condition, attraction basin, compact domain, and initial-layer assumptions are verified. It does not establish reduction of a general RFDE, preservation of global periodic-orbit folds, or infinite-horizon safety.

A nonlinear analysis must distinguish: (i) a local bifurcation in a reduced model; (ii) persistence under coupling; and (iii) intersection of an attractor, basin, or reachable tube with the unsafe set. A Hopf or fold is not automatically a sustainability transition.

\paragraph{Conjecture.} Persistence of a reduced hybrid/RFDE nonlinear transition requires a well-posed semiflow, fast difference-operator contractivity, spectral separation, center-manifold reduction, transverse Poincar\'e-map conditions, regular coupling, and preservation of material feasibility and safety. At a periodic-orbit fold, normal hyperbolicity may apply transversely but not to the fold orbit itself.

\section{Claim-status summary}
\begin{center}
\begin{tabular}{lll}
\toprule
Object & Status & Destination/scope\\
\midrule
Hybrid moiety balance & conditional theorem & represented material boundary\\
Hybrid positivity & conditional theorem & specified RFDE/hybrid solution class\\
Full-information viability & conditional benchmark & finite-dimensional ODE\\
Epistemic predecessor kernel & conditional theorem & sampled information-state system\\
Information refinement & conditional theorem & fixed authority/dynamics\\
Common-action obstruction & proposition & no observation before action\\
Observer-to-safety margin & local proposition & specified estimator error bound\\
Delay-sign statement & proposition & non-universality only\\
RFDE reduction/fold persistence & conjecture & model-specific frontier\\
\bottomrule
\end{tabular}
\end{center}

\section{Conclusion}
The paper provides a conditional mathematical spine for institution-constrained sustainability: hybrid material feasibility, information-state predecessor kernels, information monotonicity, common-action obstruction, and local observer-margin transfer. Concrete kernels, estimators, and nonlinear safety transitions remain domain-specific mathematical and empirical tasks.

\end{document}
